% % % % % % % % % % % % % % % % % % % % % % % % % % % %
%                                                     %
%   73 Suite Documentation                            %
%                                                     %
%   Author:    SCHWARZER, Andre                       %
%              von RHEIN, Christopher                 %
%              YURTBAY,   Emre                        %
%                                                     %
%   Language:  German                                 %
%                                                     %
%   Date:      22. March 2014                         %
%                                                     %
% % % % % % % % % % % % % % % % % % % % % % % % % % % %

\documentclass[a4paper,12pt]{article}

\usepackage{amssymb} % Mathe
\usepackage{amsmath} % Mathe
\usepackage[utf8x]{inputenc} % Umlaute
\usepackage[ngerman]{babel}  % Umlaute
\usepackage[T1]{fontenc}     % Umlaute im PDF
\usepackage[margin=2.5cm]{geometry} % Layout
\usepackage{booktabs}
\usepackage{times}
\usepackage{colortbl}
\usepackage{caption}
\usepackage{float}
\usepackage{multirow}
\usepackage{cite}
\usepackage{color}



% links
\usepackage[colorlinks=true, linkcolor=blue, citecolor=blue,]{hyperref}  

% glossar
\usepackage[nonumberlist]{glossaries} 

\usepackage{hyperref}

\ifx\pdftexversion\undefined
\usepackage[dvips]{graphicx}
\else
\usepackage[pdftex]{graphicx}
\DeclareGraphicsRule{*}{mps}{*}{}
\fi


\addto\captionsngerman{
 \renewcommand\glossaryname{}}
\renewcommand*{\glstextformat}[1]{\textit{\textcolor{black}{#1}} }
\makeglossary 

\newglossaryentry{CMS}
{
  name=CMS,
  description={Content Management System}
}

\newglossaryentry{ERP}
{
  name=ERP,
  description={Enterprise Resource Planning}
} % Glossar folgt


% % % % % % % % % % % % % % % % % % % % % % % % % % % %
% Variablen                                           %
% % % % % % % % % % % % % % % % % % % % % % % % % % % %
\newcommand{\authorName}{André SCHWARZER, Christopher von RHEIN, Emre YURTBAY}
\newcommand{\auftraggeber}{Hochschule Bochum - University of Applied Sciences}
\newcommand{\auftragnehmer}{Studentisches Projekt}
\newcommand{\projektName}{Suite 73 - All-In-One Data Management}
\newcommand{\documentation}{Dokumentation}
\newcommand{\tags}{\authorName}
\newcommand{\glossarName}{Erläuterung zu Begriffen und Abkürzungen}
\title{\projektName}
\author{\authorName}
\date{\today}

\definecolor{hellgrau}{rgb}{0.95,0.95,0.95}
\definecolor{koenigsblau}{rgb}{0.00, 0.09, 0.45}

% % % % % % % % % % % % % % % % % % % % % % % % % % % %
% PDF Meta Informationen                              %
% % % % % % % % % % % % % % % % % % % % % % % % % % % %
\hypersetup{
  pdfauthor   = {\authorName},
  pdfkeywords = {\tags},
  pdftitle    = {\title~(Doku)}
} 
 
% % % % % % % % % % % % % % % % % % % % % % % % % % % %
% Spezielle Formatierung für Anforderungen            %
% % % % % % % % % % % % % % % % % % % % % % % % % % % %
% \newcommand\addrow[2]{#1 &#2\\ }
% 
% \newcommand\addheading[2]{#1 &#2\\ \hline}
% \newcommand\tabularhead{\begin{tabular}{lp{13cm}}
% \hline
% }
% 
% \newcommand\addmulrow[2]{ \begin{minipage}[t][][t]{2.5cm}#1\end{minipage}% 
%    &\begin{minipage}[t][][t]{8cm}
% \begin{enumerate} #2   \end{enumerate}
%     \end{minipage}\\ }
% 
% \newenvironment{usecase}{\tabularhead}
% {\hline\end{tabular}}


\addto\captionsngerman{\renewcommand{\refname}{}}

% % % % % % % % % % % % % % % % % % % % % % % % % % % %
% Dokument                                            %
% % % % % % % % % % % % % % % % % % % % % % % % % % % %
\begin{document}
    \pagenumbering{roman}
    \begin{titlepage}
\newgeometry{left=7.5cm} 
\pagecolor{koenigsblau}
\noindent
\includegraphics[width=2cm]{73.png}\\[-1em]
\color{white}
\makebox[0pt][l]{\rule{1.3\textwidth}{1pt}}
\par
\noindent
\textbf{\textsf{Hochschule Bochum}} \textcolor{hellgrau}{\textsf{University of
Applied Sciences}}
\vfill
\noindent
{\huge \textsf{\documentation}}
\vskip\baselineskip
\noindent
\textsf{\today}
\end{titlepage}
\restoregeometry % restores the geometry
\nopagecolor% Use this to restore the color pages to white
    \setcounter{page}{2}
    \setcounter{tocdepth}{4}
    \setcounter{secnumdepth}{4}
    \tableofcontents
    \clearpage
    \pagenumbering{arabic}
 
	\begin{abstract}
		Abstrakte Beschreibung.
	\end{abstract}	
	\clearpage
 
	\section{Lastenheft}
 	
		\subsection{Einleitung}
 		Als Basis des Projekts wird eine Datenverwaltungssoftware - Base73 - entwickelt, die
eine Benutzerrechteverwaltung sowie eine allgemeine Datenverwaltung bietet.
Die Benutzerrechteverwaltung hat zwei Zuständigkeiten:
\begin{enumerate}
	\item Zugriffsrechteverwaltung auf die Verwaltungssoftware
	\item Zugriffsrechteverwaltung auf die Produkte (CMS, ERP, CRM, \ldots)
\end{enumerate}
Die Datenverwaltung ist Schnittstelle zwischen den installierbaren Systemen, wie CMS, ERP, CRM, \ldots, und der MySQL Datenbank.
 		
		\subsection{Ist-Zustand}
 		Der aktuelle Markt stellt keine einheitliche Lösung für die Verarbeitung und Verwaltung von Inhalten, Mitarbeiter, Kunden, Händler und andere, für einen Unternehmen relevante, Positionen zur Verfügung, die für ein kleines oder mittelständisches Unternehmen bezahlbar sind.\\
\\
Darüber hinaus muss man momentan aus einer Sammlung von diversen Programmen ein System selbst zusammenstellen, welches den gegebenen Anforderungen entspricht.\\
Hierbei kann es zu erhöhtem Kosten- und Zeitaufwand durch die Administration und Migration der Daten kommen.\\
Des Weiteren fehlt in den Unternehmen oft das nötige Fachwissen, um sich ein derartiges System zusammenstellen zu können.

 		
		\subsection{Soll-Konzept}
 		Beschreibung des Soll-Konzepts
	Der jetzige Markt erlaubt es Klein- sowie Mittelstandsunternehmen nicht, eine Data Management Lösung wahrnehmen zu können, da solche Programme für Großunternehmen entwickelt und nur als SAAS mit beschränkter Bandbreite angeboten werden.
	Die Suite 73 soll daher diesen Unternehmen die Möglichkeit bieten, eine All-In-One Data Management Lösung kostengünstig nutzen zu können.
	Die Base 73 soll ein CMS enthalten und alle Funktionen frei anbieten. Erweiterung von CRM, ERP bis hin zum Onlineshop sollen optional kostengünstig zu erwerben sein.
	Zu den Inhalten sollen SEO & SEM's gegeben sein, ein Frontend und ein Backend  mit den entsprechenden Funktionen (Drag & Drop)sollten auch bereit stehen.
 	\clearpage

	\section{Pflichten}
 	Hier die Pflichen

\end{document}

